% =====================================================
% settings.tex
% Cấu hình chung cho báo cáo
% Sinh viên chủ yếu sửa thông tin cá nhân ở mục 7.
% =====================================================

% =========================
% 1. CẤU HÌNH TIẾNG VIỆT VÀ FONT
% =========================
\usepackage{fontspec}
\IfFontExistsTF{Times New Roman}
  {\setmainfont{Times New Roman}}
  {\setmainfont{Tinos}}

\usepackage{polyglossia}
\setdefaultlanguage{vietnamese}

% =========================
% 2. CẤU HÌNH TRANG
% =========================
\usepackage[
    a4paper,
    left=3.0cm,
    right=1.5cm,
    top=2.0cm,
    bottom=2.0cm
]{geometry}

\usepackage{setspace}
\onehalfspacing

% Cỡ chữ mặc định của nội dung là 13pt. Các phần có quy định riêng
% bằng \fontsize trong tài liệu vẫn giữ nguyên cỡ chữ được chỉ định.
\makeatletter
\renewcommand{\normalsize}{%
    \@setfontsize\normalsize{13pt}{15.6pt}%
    \abovedisplayskip 13pt plus 3pt minus 7pt%
    \abovedisplayshortskip 0pt plus 4pt%
    \belowdisplayshortskip 6.5pt plus 3.5pt minus 3pt%
    \belowdisplayskip \abovedisplayskip
    \let\@listi\@listI
}
\makeatother
\normalsize

\usepackage{indentfirst}
\setlength{\parindent}{1.27cm}
\setlength{\parskip}{6pt}

% =========================
% 3. GÓI HỖ TRỢ
% =========================
\usepackage{graphicx}
\usepackage{float}
\usepackage{array}
\usepackage{longtable}
\usepackage{booktabs}
\usepackage{multirow}
\usepackage{tabularx}
\usepackage{caption}
\usepackage{amsmath, amssymb}
\usepackage{enumitem}
\setlist[itemize]{
    label=-,
    leftmargin=1.27cm,
    labelsep=0.5em
}
\usepackage{titlesec}
\usepackage{tocloft}
\usepackage{hyperref}
\usepackage{tikz}
\usepackage{eso-pic}
\usepackage{fancyhdr}

% =========================
% 4. HYPERLINK
% =========================
\hypersetup{
    colorlinks=true,
    linkcolor=black,
    citecolor=black,
    urlcolor=blue
}

% Số trang đặt giữa phía trên; áp dụng cả cho trang mở đầu chương.
\setlength{\headheight}{17pt}
\pagestyle{fancy}
\fancyhf{}
\fancyhead[C]{\thepage}
\renewcommand{\headrulewidth}{0pt}
\fancypagestyle{plain}{
    \fancyhf{}
    \fancyhead[C]{\thepage}
    \renewcommand{\headrulewidth}{0pt}
}

% =========================
% 5. ĐÁNH SỐ HÌNH, BẢNG THEO CHƯƠNG
% =========================
\counterwithin{table}{chapter}
\counterwithin{figure}{chapter}
\renewcommand{\thetable}{\thechapter.\arabic{table}}
\renewcommand{\thefigure}{\thechapter.\arabic{figure}}

% Môi trường riêng cho sơ đồ để đánh số và lập danh mục độc lập.
\newfloat{sodo}{H}{los}[chapter]
\floatname{sodo}{Sơ đồ}
\renewcommand{\thesodo}{\thechapter.\arabic{sodo}}
\makeatletter
\newcommand{\l@sodo}[2]{%
    \begingroup
    \let\ReportOriginalNumberline\numberline
    \renewcommand{\numberline}[1]{%
        \ReportOriginalNumberline{Sơ đồ~##1}}%
    \@dottedtocline{1}{0pt}{2.8cm}{#1}{#2}%
    \endgroup
}
\makeatother

\captionsetup{
    font=normalsize,
    labelfont=normalfont,
    textfont=normalfont,
    labelsep=period,
    justification=centering
}
\captionsetup[table]{name=Bảng,position=top}
\captionsetup[figure]{name=Hình,position=bottom}
\captionsetup[sodo]{name=Sơ đồ,position=bottom}

\renewcommand{\listtablename}{BẢNG}
\renewcommand{\listfigurename}{HÌNH, SƠ ĐỒ}
\renewcommand{\contentsname}{MỤC LỤC}
\renewcommand{\bibname}{TÀI LIỆU THAM KHẢO}

% =========================
% 6. ĐỊNH DẠNG TIÊU ĐỀ CHƯƠNG, MỤC
% =========================
\titleformat{\chapter}
  [display]
  {\centering\bfseries\fontsize{14pt}{16pt}\selectfont}
  {Chương \thechapter}
  {0.3cm}
  {\MakeUppercase}

\titlespacing*{\chapter}{0pt}{0pt}{18pt}

\titleformat{\section}
  {\bfseries\fontsize{13pt}{15pt}\selectfont}
  {\thesection}
  {0.5em}
  {}

\titleformat{\subsection}
  {\bfseries\itshape\fontsize{13pt}{15pt}\selectfont}
  {\thesubsection}
  {0.5em}
  {}

\titleformat{\subsubsection}
  {\itshape\fontsize{13pt}{15pt}\selectfont}
  {\thesubsubsection}
  {0.5em}
  {}

\setcounter{secnumdepth}{3}
\setcounter{tocdepth}{3}

% =========================
% 7. ĐỊNH DẠNG MỤC LỤC
% =========================
\renewcommand{\cftchapfont}{\bfseries\fontsize{12pt}{14pt}\selectfont}
\renewcommand{\cftchappagefont}{\bfseries\fontsize{12pt}{14pt}\selectfont}
\renewcommand{\cftchappresnum}{CHƯƠNG~}
\renewcommand{\cftchapaftersnum}{.}
\setlength{\cftchapnumwidth}{2.5cm}
\renewcommand{\cftsecfont}{\bfseries\fontsize{12pt}{14pt}\selectfont}
\renewcommand{\cftsecpagefont}{\bfseries\fontsize{12pt}{14pt}\selectfont}
\renewcommand{\cftsecaftersnum}{.}
\renewcommand{\cftsubsecfont}{\itshape\fontsize{12pt}{14pt}\selectfont}
\renewcommand{\cftsubsecpagefont}{\fontsize{12pt}{14pt}\selectfont}
\renewcommand{\cftsubsecaftersnum}{.}
\renewcommand{\cftsubsubsecfont}{\fontsize{12pt}{14pt}\selectfont}
\renewcommand{\cftsubsubsecpagefont}{\fontsize{12pt}{14pt}\selectfont}
\renewcommand{\cftsubsubsecaftersnum}{.}
\renewcommand{\cftchapleader}{\cftdotfill{\cftdotsep}}
\setlength{\cftbeforechapskip}{6pt}

\renewcommand{\cfttoctitlefont}{\hfill\bfseries\fontsize{14pt}{16pt}\selectfont\MakeUppercase}
\renewcommand{\cftaftertoctitle}{\hfill}
\setlength{\cftbeforetoctitleskip}{0pt}
\setlength{\cftaftertoctitleskip}{18pt}

\makeatletter
\newcommand{\ReportListTable}[4]{%
    {\bfseries\fontsize{13pt}{14pt}\selectfont #1:\par}
    \vspace{6pt}
    \begingroup
    \fontsize{12pt}{13pt}\selectfont
    \renewcommand{\arraystretch}{1.5}
    \def\ReportListPrefix{#3}%
    \def\numberline##1{\ReportListPrefix~##1.\enspace}%
    \def\contentsline##1##2##3##4{##2\dotfill & ##3\\\hline}%
    \def\addvspace##1{}%
    \def\select@language##1{}%
    \makeatletter
    \noindent
    \begin{tabular}{|>{\raggedright\arraybackslash}p{\dimexpr\textwidth-3.5cm\relax}|>{\centering\arraybackslash}p{2.6cm}|}
        \hline
        \multicolumn{1}{|c|}{\bfseries #2} & \bfseries Trang \\
        \hline
        \@input{\jobname.#4}%
    \end{tabular}
    \par
    % Mở lại tệp danh mục để các caption phía sau tiếp tục ghi dữ liệu.
    \let\contentsline\@gobblefour
    \let\addvspace\@gobble
    \let\select@language\@gobble
    \@starttoc{#4}%
    \makeatother
    \endgroup
}

\newcommand{\CombinedLists}{%
    \chapter*{\fontsize{13pt}{14pt}\selectfont DANH MỤC CÁC BẢNG, SƠ ĐỒ, HÌNH}
    \addcontentsline{toc}{chapter}{DANH MỤC CÁC BẢNG, SƠ ĐỒ, HÌNH}
    \ReportListTable{BẢNG}{Tên bảng}{Bảng}{lot}
    \vspace{14pt}
    \ReportListTable{SƠ ĐỒ}{Tên sơ đồ}{Sơ đồ}{los}
    \vspace{14pt}
    \ReportListTable{HÌNH}{Tên hình}{Hình}{lof}
}
\makeatother

% =========================
% 8. THÔNG TIN SINH VIÊN/BÁO CÁO
% Sinh viên chỉ cần sửa phần này.
% =========================
\newcommand{\BoGiaoDuc}{BỘ GIÁO DỤC VÀ ĐÀO TẠO}
\newcommand{\TenTruong}{TRƯỜNG ĐẠI HỌC ĐẠI NAM}

\newcommand{\TenSinhVien}{ĐINH VĂN TÂN LƯỢNG}
\newcommand{\MaSinhVien}{1771020446}
\newcommand{\KhoaHoc}{17}

\newcommand{\TenBaoCao}{TÊN BÁO CÁO}
\newcommand{\LoaiBaoCao}{BÁO CÁO THỰC TẬP TỐT NGHIỆP CHUYÊN NGÀNH: CÔNG NGHỆ THÔNG TIN}
\newcommand{\Nganh}{CÔNG NGHỆ THÔNG TIN}
\newcommand{\ChuyenNganh}{LẬP TRÌNH NHÚNG VÀ IOT}

\newcommand{\GiangVienHD}{ThS. Nguyễn Văn Nhân}

\newcommand{\DiaDiem}{HÀ NỘI}
\newcommand{\Nam}{2026}

% Logo trường. Có thể thay bằng file khác trong thư mục images/.
\newcommand{\LogoTruong}{images/logo-dai-nam.png}

% =========================
% 9. LỆNH CHƯƠNG KHÔNG ĐÁNH SỐ
% Dùng cho: Lời cam đoan, Lời cảm ơn, Mở đầu, Kết luận...
% =========================
\newcommand{\UnnumberedChapter}[1]{
    \chapter*{#1}
    \addcontentsline{toc}{chapter}{#1}
}

% Ghi nguồn ngay dưới bảng, hình hoặc sơ đồ khi sử dụng dữ liệu bên ngoài.
\newcommand{\Nguon}[1]{%
    \par\smallskip
    {\centering\itshape Nguồn: #1\par}
}

% Với môi trường sodo của gói float, caption luôn được dồn xuống cuối.
% Gộp nguồn vào caption để bảo đảm thứ tự hiển thị: tên sơ đồ, rồi đến nguồn.
\newcommand{\SodoCaption}[2]{%
    \caption[#1]{#1\\[3pt]\textit{Nguồn: #2}}%
}

% =========================
% 10. KHUNG TRANG BÌA
% =========================
\newcommand{\DrawCoverBorder}{
\AddToShipoutPictureBG*{
\begin{tikzpicture}[remember picture, overlay]
    \draw[line width=1pt]
        ([xshift=1.2cm,yshift=-1.2cm]current page.north west)
        rectangle
        ([xshift=-1.2cm,yshift=1.2cm]current page.south east);
\end{tikzpicture}
}
}
